\section{Methodology}

We compare a direct PINN, which minimizes a penalized form of the constrained objective, with an indirect PINN, which solves the first-order optimality system. Both methods use a single control network $U_\psi(x,t)$. To keep the residual notation common to both formulations, define the state-equation residual operator
\begin{equation}
\operatorname{Res}_y[v]:=\partial_tv-\nu\Delta v+f(v)-g.
\label{eq:state_residual_operator}
\end{equation}

\subsection{Direct PINN Formulation}

The state and control are approximated by $Y_\theta(x,t)$ and $U_\psi(x,t)$. Their state-equation residual is $\mathcal R_y:=\operatorname{Res}_y[Y_\theta]$.
The direct loss is
\begin{equation}
\label{eq:direct_loss}
\mathcal L_{\rm dir}
=w_{\rm res}\mathcal L_{\rm res}^y
+w_{\rm bc}\mathcal L_{\rm bc}^y
+w_{\rm ic}\mathcal L_{\rm ic}^y
+\frac12\mathcal L_{\rm track}
+\frac{\alpha}{2}\mathcal L_{\rm reg},
\end{equation}
where
\begin{align}
\mathcal L_{\rm res}^y&:=\frac1{N_{\rm int}}\sum_{j=1}^{N_{\rm int}}|\mathcal R_y(x_j,t_j)|^2,\\
\mathcal L_{\rm bc}^y&:=\frac1{N_{\rm bc}}\sum_{j=1}^{N_{\rm bc}}|Y_\theta(x_j^b,t_j^b)-U_\psi(x_j^b,t_j^b)|^2,\\
\mathcal L_{\rm ic}^y&:=\frac1{N_{\rm ic}}\sum_{i=1}^{N_{\rm ic}}|Y_\theta(x_i,0)-y_0(x_i)|^2,\\
\mathcal L_{\rm track}&:=\frac1{N_{\rm int}}\sum_{j=1}^{N_{\rm int}}|Y_\theta(x_j,t_j)-y_d(x_j,t_j)|^2.
\end{align}
In one space dimension the boundary consists of two endpoints, and the control regularization uses counting measure on $\partial\Omega$:
\begin{equation}
\mathcal L_{\rm reg}:=\frac1{N_{\rm bc}}\sum_{j=1}^{N_{\rm bc}}
\left(|U_\psi(0,t_j^b)|^2+|U_\psi(1,t_j^b)|^2\right).
\end{equation}
This formulation is simple and avoids an adjoint network, but boundary satisfaction depends on the penalty weight and can be delicate~\cite{zhang2026pinns}.

\subsection{Indirect PINN Formulation with Hard Constraints}

\subsubsection{First-order optimality system}

With
\begin{equation}
\mathcal{L}(y,u,\lambda)=J(y,u)+\int_Q\lambda\left(\partial_ty-\nu\Delta y+f(y)-g\right)\,dx\,dt,
\end{equation}
integration by parts gives
\begin{align}
\partial_ty-\nu\Delta y+f(y)-g&=0 &&\text{in }Q, \label{eq:state_opt}\\
-\partial_t\lambda-\nu\Delta\lambda+f'(y)\lambda+(y-y_d)&=0 &&\text{in }Q, \label{eq:adjoint}\\
\alpha u+\nu\partial_n\lambda&=0 &&\text{on }\Sigma, \label{eq:stationarity}\\
y=u,\qquad \lambda&=0 &&\text{on }\Sigma, \label{eq:bc_y}\\
y(x,0)=y_0(x),\qquad \lambda(x,T)&=0 &&\text{in }\Omega. \label{eq:terminal}
\end{align}
For $\Omega=(0,1)$, $\partial_n\lambda(0,t)=-\lambda_x(0,t)$ and $\partial_n\lambda(1,t)=\lambda_x(1,t)$.

\subsubsection{One control network and hard constraints}

We use one network $U_\psi(x,t)$ and define its two boundary traces by
\begin{equation}
u_0(t):=U_\psi(0,t),\qquad u_1(t):=U_\psi(1,t).
\end{equation}
For $\Omega=(0,1)$, the boundary lifting and state trial function are
\begin{align}
B_\psi(x,t)&:=(1-x)U_\psi(0,t)+xU_\psi(1,t),\\
Y_\theta(x,t)&:=B_\psi(x,t)+x(1-x)N_\theta(x,t). \label{eq:hard_state}
\end{align}
Thus $Y_\theta(0,t)=U_\psi(0,t)$ and $Y_\theta(1,t)=U_\psi(1,t)$ for every parameter value. The adjoint is parameterized as
\begin{equation}
\Lambda_\phi(x,t):=x(1-x)(T-t)M_\phi(x,t), \label{eq:hard_adjoint}
\end{equation}
which enforces both $\Lambda_\phi=0$ on the spatial boundary and $\Lambda_\phi(x,T)=0$. The state initial condition remains a soft constraint.

\subsubsection{Indirect PINN loss}

Define
\begin{align}
\mathcal R_y&:=\operatorname{Res}_y[Y_\theta], \label{eq:res_y}\\
\mathcal R_\lambda&:=-\partial_t\Lambda_\phi-\nu\Delta\Lambda_\phi+f'(Y_\theta)\Lambda_\phi+Y_\theta-y_d, \label{eq:res_lambda}\\
\mathcal R_u&:=\alpha U_\psi+\nu\partial_n\Lambda_\phi &&\text{on }\Sigma. \label{eq:res_u}
\end{align}
In one space dimension, its two traces are
\begin{equation}
\mathcal R_{u,0}(t):=\mathcal R_u(0,t)=\alpha U_\psi(0,t)-\nu\partial_x\Lambda_\phi(0,t),\qquad
\mathcal R_{u,1}(t):=\mathcal R_u(1,t)=\alpha U_\psi(1,t)+\nu\partial_x\Lambda_\phi(1,t).
\label{eq:stationarity_traces}
\end{equation}
The indirect loss is
\begin{equation}
\label{eq:indirect_loss}
\mathcal L_{\rm ind}=w_y\mathcal L_y+w_\lambda\mathcal L_\lambda+w_u\mathcal L_u+w_{\rm ic}\mathcal L_{\rm ic},
\end{equation}
with
\begin{align}
\mathcal L_y&:=\frac1{N_{\rm int}}\sum_j|\mathcal R_y(x_j,t_j)|^2,\\
\mathcal L_\lambda&:=\frac1{N_{\rm int}}\sum_j|\mathcal R_\lambda(x_j,t_j)|^2,\\
\mathcal L_u&:=\frac1{2N_{\rm bc}}\sum_j\left(|\mathcal R_{u,0}(t_j^b)|^2+|\mathcal R_{u,1}(t_j^b)|^2\right), \label{eq:boundary_residual_loss}\\
\mathcal L_{\rm ic}&:=\frac1{N_{\rm ic}}\sum_i|Y_\theta(x_i,0)-y_0(x_i)|^2.
\end{align}
Thus $\mathcal L_u$ is the sample mean with respect to time and the normalized counting measure on the two endpoints; the factor $1/2$ is part of its definition throughout the paper.
There are no state-boundary, adjoint-boundary, or adjoint-terminal penalties because these constraints are satisfied by construction.
